\begin{table}[H]
\centering
\caption{Standardized Effect Sizes for Main Outcomes}
\label{tab:sde}
\begin{adjustbox}{max width=\textwidth}
\begin{tabular}{lcccccc}
\toprule
Outcome & $\hat{\beta}$ & SE & SD($Y$) & SDE & SE(SDE) & Classification \\
\midrule
\multicolumn{7}{l}{\textit{Panel A: Pooled}} \\
log(NWP) & 0.0699 & (0.1146) & 1.130 & 0.0619 & (0.1014) & Moderate positive \\
Loss Ratio (\%) & 6.2881 & (5.7499) & 14.549 & 0.4322 & (0.3952) & Large positive \\
\midrule
\multicolumn{7}{l}{\textit{Panel B: Heterogeneous (Motor vs.~Property)}} \\
log(NWP) --- Motor & 0.0084 & (0.1171) & 1.130 & 0.0074 & (0.1037) & Small positive \\
log(NWP) --- Property & 0.1929 & (0.0850) & 1.130 & 0.1707 & (0.0752) & Large positive \\
Loss Ratio --- Motor & 11.9544 & (2.7284) & 14.549 & 0.8217 & (0.1875) & Large positive \\
Loss Ratio --- Property & -5.0444 & (2.6383) & 14.549 & -0.3467 & (0.1813) & Large negative \\
\bottomrule
\end{tabular}
\end{adjustbox}

\vspace{0.3em}
\begin{minipage}{\textwidth}
\footnotesize
\textit{Notes:} \textbf{Country:} United Kingdom. \textbf{Research question:} Did the FCA's ban on insurance price-walking (GIPP, January 2022) create a waterbed effect in motor and home insurance, where new customer premiums rose to compensate for the loss of renewal overcharges? \textbf{Policy mechanism:} GIPP (ICOBS 6B) requires that renewal prices must not exceed the Equivalent New Business Price, eliminating the practice of charging loyal customers higher premiums than new customers in home and motor insurance markets. \textbf{Outcome definition:} Log Net Written Premium (quarterly, by line of business) and Loss Ratio (claims incurred as percentage of premiums earned) from BoE data; Claims Frequency and Complaints Rate (band midpoints) from FCA GI Value Measures. \textbf{Treatment:} Binary; GIPP-targeted lines (Motor liability, Motor other, Property) vs.\ non-targeted lines (8 other non-life classes). \textbf{Data:} Bank of England Insurance Aggregate Data (2017Q1--2025Q4, 396 line-quarter obs) and FCA GI Value Measures (2021--2024, 1,636 firm-product-year obs). \textbf{Method:} Two-way fixed effects DiD (line + quarter FE), clustered at line level for BoE and firm level for FCA data. \textbf{Sample:} 11 non-life insurance lines over 36 quarters (BoE); 184 firms across 55 product categories over 3 years (FCA). SDE $= \hat{\beta} / \text{SD}(Y)$ where SD($Y$) is the unconditional standard deviation. Classification refers to magnitude, not statistical significance: Large ($|$SDE$| > 0.15$), Moderate ($0.05$--$0.15$), Small ($0.005$--$0.05$), Null ($< 0.005$).
\end{minipage}
\end{table}
